% SPDX-License-Identifier: Apache-2.0
\section{A Pipeline Is an Async Function}
\label{sec:e-pipeline}

A pipeline and a sequence are one construct driven two ways.

A sequence runs one invocation to completion before starting the next. A
pipeline starts a new invocation every cycle, so several are in flight at
once, each at a different point in the body. Nothing about the body
differs. The difference is the rate at which invocations start, and that
belongs at the call site.

\begin{lstlisting}
async fn mac(a: U<32>, b: U<32>, prev: U<64>) -> U<64> {
    let p = mul(a, b);
    tick().await;              // stage boundary
    prev + p
}
\end{lstlisting}

Three constructs go, and one of them takes a compiler pass with it.

\subsection{The pipe declaration disappears}

LHDL's \code{pipe} declared a value that crosses a stage boundary, so that
live-range analysis would know to hold it. TxHDL had the same notion under
a different name.

In an async function, a local that is live across an \code{.await} is held
in the state machine the compiler generates, because that is what an async
function is. \code{p} in the listing above crosses \code{tick().await} and
needs no declaration. The analysis the \code{pipe} keyword existed to feed
is one that \code{rustc} performs already.

\subsection{What this does not settle}

Three things, and none of them is free.

\begin{itemize}
  \item \term{Named stages are gone.} A boundary is an await, not a named
        block. A name becomes a comment.
  \item \term{The initiation interval moves to the call site.} It was
        implied by the choice of keyword and is now an argument where the
        body is driven. That is a fair trade and not a saving.
  \item \term{Resource sharing has to be proved.} A real pipeline shares
        one multiplier across every item in flight, while several
        concurrent invocations each appear to want their own. They are the
        same code at the same await index, so the lowering can share them.
        That is a claim about the compiler being written, not about Rust,
        and it is the largest piece of unfinished work in this proposal.
\end{itemize}

One thing improves. TxHDL had a \code{stall} statement for a stage that
cannot proceed. A stall is now an await that has not completed, which is
the mechanism rather than a construct layered on top of it.
